\section{Summary of Contributions}
\label{sec:conclusion-summary}

This thesis set out to answer a specific research question, posed directly by the supervisor: can a
query recommender that conditions on both a user's stated queries and their real interaction
behavior produce measurably better follow-up suggestions than one conditioning on query history
alone? The answer, established in Chapter~\ref{ch:eval}, is yes: a controlled ablation, run on the
real production recommender and scored against real held-out ground truth, shows a $6.5\times$
increase in behavioral-alignment@3 and a more than $2\times$ increase in proxy-recall@3 when the
behavioral signal is included, consistently across every evaluated user and question.

Reaching that result required, in order: a per-user behavioral divergence signal engineered from two
data sources that share no natural key (Chapter~\ref{ch:behavior}); a retrieval-augmented,
LLM-based candidate-generation step conditioned on that signal (Section~\ref{sec:rec-generation});
and a ranking stage that makes the signal load-bearing in the final output
(Section~\ref{sec:rec-ranking}). One design goal, however, was pursued and found not to be
achievable within this thesis's scope: a precision schema-grounding filter, motivated by a
well-documented hallucination risk in the LLM+KG literature, was calibrated against this project's
actual ontologies and found unable to separate the supervisor's own canonical example query from an
off-topic control phrase, because the underlying schema metadata is too sparse for embedding
similarity to carry the necessary signal. This is reported as a concrete, empirically-grounded
negative result (Section~\ref{sec:rec-validation}) rather than an over-claimed feature.

\section{Limitations}
\label{sec:conclusion-limits}

Three limitations bound how far this thesis's results can be claimed to generalize, beyond those
already discussed in each chapter. First, all quantitative results are obtained on a single,
synthetic, Faker-generated dataset with no true causal link between stated search intent and later
watch behavior; the \emph{direction} of the observed divergence and ablation effects is consistent
with the supervisor's own real-world framing, but their \emph{magnitude} is, in part, a property of
this dataset's generation process. Second, per the scope decision in
Section~\ref{sec:arch-vkgs}, the entire behavioral, stated-vs-actual contribution is demonstrable
only on the Movie domain; the Tourism domain's only large behavioral signal is anonymous and was
deliberately excluded. Third, the recommender's input-contract dependency on an external
application --- what real interaction data a caller actually logs per user --- was never resolved
during this thesis; the \code{interaction\_data} field in the API contract
(Section~\ref{sec:arch-api}) remains an untyped placeholder, and all behavioral conditioning in this
work instead uses the offline, precomputed profile described in Chapter~\ref{ch:behavior}.

\section{Future Work}
\label{sec:conclusion-future}

Several concrete extensions follow directly from limitations identified in this thesis, rather than
being open-ended:

\begin{itemize}
    \item \textbf{A larger evaluation run.} Chapter~\ref{ch:eval}'s ablation used 3 users and 4
        questions by explicit scope choice; the evaluation harness already supports a larger run via
        configuration constants alone, with no design or code changes required.
    \item \textbf{Validation against real interaction data.} If the external application's logging
        contract is ever confirmed (Section~\ref{sec:conclusion-limits}), or if any other source of
        real user interaction data becomes available, both the behavioral-profiling pipeline
        (Chapter~\ref{ch:behavior}) and the evaluation harness (Chapter~\ref{ch:eval}) can be
        re-run against it directly, since both already consume this exact CSV schema; this would
        directly test whether the divergence effect's magnitude, not just its direction, holds
        outside synthetic data.
    \item \textbf{Enriching schema metadata for a real precision grounding filter.} The negative
        result in Section~\ref{sec:rec-validation} traces to sparse class descriptions in the source
        ontologies; enriching them with real catalog/instance values pulled from the live VKG
        endpoints is a concrete, scoped path toward the precision hallucination filter originally
        intended.
    \item \textbf{A sensitivity analysis on heuristic parameters.} The temporal-join window
        (Section~\ref{sec:behavior-link}) and the candidate-pool size
        (Section~\ref{sec:rec-generation}) are both fixed, undertuned choices; a systematic sweep
        over both was left for future work once real interaction data makes such tuning meaningful.
    \item \textbf{Re-evaluating against a funded cloud LLM.} All generation and evaluation in this
        thesis used a local Ollama model after both configured cloud providers were blocked by
        account issues unrelated to the code; re-running with a stronger, funded model is a direct,
        no-code-change comparison worth pursuing.
\end{itemize}
